\chapter{Ringen en rekenkunde}\label{ch:b3:rings}

Gewone gehele getallen ontbinden op precies één manier in
priemfactoren; veeltermen over een lichaam ook. Zijn die twee
feiten één stelling? Dit hoofdstuk antwoordt bevestigend en legt
precies de hypothesen bloot die een ``rekenkunde'' in een
commutatieve ring mogelijk maken: de keten
\[
\text{euclidisch} \;\Longrightarrow\; \text{hoofdideaal}
\;\Longrightarrow\; \text{factorieel (UFD)},
\]
met alle implicaties bewezen en alle omkeringen weerlegd. Daarna
zetten we de theorie in waar ze haar waarde bewijst: bij de gehele
getallen van Gauss $\Z[\iu]$ (die in de weekendopgave de
tweekwadratenstelling van Fermat zullen kraken), bij
veeltermringen in meer veranderlijken (het lemma van Gauss, het
criterium van \omterm{thm:b3:rings:criteria}{Eisenstein}) en bij \omterm{def:b3:rings:noetherian}{noetherse ringen}, met als
hoogtepunt de basisstelling van Hilbert. Overal betekent
\emph{ring} een commutatieve ring met eenheid $1 \neq 0$; de
\omterm{def:b3:rings:ideal}{idealen} van $\Z$ en $K[X]$ uit het volume van bachelorjaar 2 zijn
onze twee leidende voorbeelden.

\section{Idealen, quotiënten en de isomorfiestelling}

\begin{definition}\label{def:b3:rings:ideal}
Een \emph{ideaal}\index{ideaal} $I$ van een ring $A$ is een
additieve deelgroep met $AI \subseteq I$. De
\emph{quotiëntring}\index{quotiëntring} $A/I$ is de \omterm{thm:b3:groups:quotient}{quotiëntgroep}
$(A, +)/I$ met de vermenigvuldiging $(a + I)(b + I) = ab + I$: die
is goed gedefinieerd, want $a$ vervangen door $a + x$ en $b$ door
$b + y$ (met $x, y \in I$) verandert $ab$ met $ay + xb + xy \in I$.
De projectie $\pi \colon A \to A/I$ is een surjectief ringmorfisme
met kern $I$, en de kernen van ringmorfismen zijn precies de
idealen.
\end{definition}

\begin{theorem}[Eerste isomorfiestelling]\label{thm:b3:rings:firstiso}
Is $f \colon A \to B$ een ringmorfisme, dan is $\bar f\colon A/\ker
f \to \operatorname{im} f$, $a + \ker f \mapsto f(a)$, een
ringisomorfisme. Algemener factoriseert $f$ over $A/I$ voor elk
\omterm{def:b3:rings:ideal}{ideaal} $I \subseteq \ker f$. De \omterm{def:b3:rings:ideal}{idealen} van $A/I$ zijn de $J/I$ met
$J \supseteq I$ een \omterm{def:b3:rings:ideal}{ideaal} van $A$ (correspondentiestelling).
\end{theorem}

\begin{proof}
Net als bij groepen (\cref{thm:b3:groups:firstiso,%
thm:b3:groups:correspondence}), waarbij we opmerken dat alle
betrokken afbeeldingen ook de producten respecteren: $\bar f$ is
goed gedefinieerd, bijectief op het beeld en multiplicatief; de
correspondentie $J \mapsto J/I$, $\bar J \mapsto \pi^{-1}(\bar J)$
bewaart in beide richtingen \omterm{def:b3:rings:ideal}{idealen}, omdat $\pi$ een surjectief
ringmorfisme is.
\end{proof}

\begin{definition}\label{def:b3:rings:primemaximal}
Zij $I \subsetneq A$ een echt \omterm{def:b3:rings:ideal}{ideaal}. $I$ heet
\emph{priem}\index{priemideaal} als uit $ab \in I$ volgt dat $a \in
I$ of $b \in I$; en $I$ heet \emph{maximaal}\index{maximaal
ideaal} als er geen \omterm{def:b3:rings:ideal}{ideaal} strikt tussen $I$ en $A$ ligt.
\end{definition}

\begin{proposition}\label{prop:b3:rings:primemaximal}
$I$ is priem $\iff$ $A/I$ is een integriteitsdomein; $I$ is
\omterm{def:b3:rings:primemaximal}{maximaal} $\iff$ $A/I$ is een lichaam. In het bijzonder zijn
maximale \omterm{def:b3:rings:ideal}{idealen} priem.
\end{proposition}

\begin{proof}
Schrijf $\bar a$ voor de klassen in $A/I$. ``$I$ priem'' vertaalt
zich letterlijk naar ``$\bar a\bar b = 0 \Rightarrow \bar a = 0$ of
$\bar b = 0$'', en $A/I \neq 0$ naar $I \neq A$: dat is precies de
definitie van een domein. Voor de maximaliteit gebruiken we de
correspondentiestelling: geen \omterm{def:b3:rings:ideal}{ideaal} strikt tussen $I$ en $A$
$\iff$ $A/I$ heeft geen ander \omterm{def:b3:rings:ideal}{ideaal} dan $0$ en zichzelf $\iff$
$A/I$ is een lichaam --- voor die laatste stap: in een lichaam zijn
$0$ en het geheel de enige \omterm{def:b3:rings:ideal}{idealen} (een \omterm{def:b3:rings:ideal}{ideaal} dat $x \ne 0$ bevat,
bevat $x^{-1}x = 1$); en omgekeerd, brengt elke $x \neq 0$ het
eenheidsideaal voort, dan is $xy = 1$ voor zekere $y$. Lichamen
zijn domeinen, dus zijn maximale \omterm{def:b3:rings:ideal}{idealen} priem.
\end{proof}

\begin{example}\label{ex:b3:rings:primemaximal}
In $\Z$ zijn de priemidealen $(0)$ en de $(p)$ met $p$ priem; de
maximale zijn de $(p)$ ($\Z/p\Z = \mathbb F_p$ is een lichaam,
$\Z/(0) = \Z$ niet). In $K[X, Y]$ zijn $(X) \subsetneq (X, Y)$
beide priem ($K[X,Y]/(X) \cong K[Y]$, een domein; $K[X,Y]/(X,Y)
\cong K$, een lichaam), zodat $(X)$ priem maar niet \omterm{def:b3:rings:primemaximal}{maximaal} is.
\end{example}

Om in volle algemeenheid te garanderen dat er maximale \omterm{def:b3:rings:ideal}{idealen}
\emph{bestaan}, hebben we een verzamelingtheoretisch principe
nodig. Een partieel geordende verzameling heet \emph{inductief} als
elke totaal geordende deelverzameling (\emph{keten}) een
bovengrens heeft.

\begin{theorem}[Lemma van Zorn]\label{thm:b3:rings:zorn}
Elke niet-lege inductief geordende verzameling bezit een \omterm{def:b3:rings:primemaximal}{maximaal}
element.\index{lemma van Zorn}
\admitted
\end{theorem}

\begin{remark}\label{rem:b3:rings:zorn}
Dit is geen stelling van de gewone wiskunde maar een
\emph{axioma}: het is, boven op de basisaxioma's van
Zermelo--Fraenkel voor de verzamelingenleer, gelijkwaardig met het
keuzeaxioma (``elk product van niet-lege verzamelingen is
niet-leeg''), dat we in dit boek overal aanvaarden. We signaleren
elk gebruik ervan. De analyse zal het opnieuw inroepen
(Hahn--Banach, \cref{ch:b3:banach}).
\end{remark}

\begin{theorem}[Krull]\label{thm:b3:rings:krull}
Elk echt \omterm{def:b3:rings:ideal}{ideaal} $I \subsetneq A$ ligt in een \omterm{def:b3:rings:primemaximal}{maximaal
ideaal}.\index{stelling van Krull}
\end{theorem}

\begin{proof}
Orden de verzameling $\mathcal E$ van echte \omterm{def:b3:rings:ideal}{idealen} die $I$
bevatten naar inclusie; ze is niet leeg ($I \in \mathcal E$). Een
keten $(J_\lambda)$ in $\mathcal E$ heeft als bovengrens $J =
\bigcup J_\lambda$: dat is een \omterm{def:b3:rings:ideal}{ideaal} (twee elementen $a, b \in J$
liggen wegens totaliteit in een gemeenschappelijke $J_\lambda$),
het is echt ($1 \notin J_\lambda$ voor alle $\lambda$) en het bevat
$I$. Het lemma van Zorn levert een \omterm{def:b3:rings:primemaximal}{maximaal} element van $\mathcal
E$, en dat is een \omterm{def:b3:rings:primemaximal}{maximaal ideaal} dat $I$ bevat (een echt \omterm{def:b3:rings:ideal}{ideaal}
strikt daarboven zou in $\mathcal E$ liggen).
\end{proof}

\begin{theorem}[Chinese reststelling]\label{thm:b3:rings:crt}
Zij $I_1, \dots, I_n$ paarsgewijs
\emph{comaximale}\index{comaximale idealen} \omterm{def:b3:rings:ideal}{idealen} van $A$ ($I_k +
I_l = A$ voor $k \neq l$). Dan is
\[
A\Big/\bigcap_{k=1}^n I_k \;\xrightarrow{\;\sim\;}\;
\prod_{k=1}^n A/I_k,
\qquad
a \longmapsto (a + I_1, \dots, a + I_n),
\]
en bovendien $\bigcap_k I_k = I_1 I_2 \cdots I_n$ (het door de
producten voortgebrachte \omterm{def:b3:rings:ideal}{ideaal}).\index{Chinese reststelling}
\end{theorem}

\begin{proof}
De afbeelding $f(a) = (a + I_k)_k$ is een ringmorfisme met kern
$\bigcap I_k$; wegens \cref{thm:b3:rings:firstiso} volstaat het de
surjectiviteit te bewijzen. Leg $k$ vast; schrijf voor elke $l \neq
k$ de gelijkheid $1 = u_l + v_l$ met $u_l \in I_k$ en $v_l \in I_l$
(comaximaliteit). Dan is
\[
e_k = \prod_{l \neq k} v_l = \prod_{l\neq k}(1 - u_l)
\equiv 1 \pmod{I_k},
\qquad e_k \in I_l \ (l \neq k),
\]
dus $f(e_k) = (0, \dots, 1, \dots, 0)$; is een doel $(a_k + I_k)_k$
gegeven, dan wordt het bereikt door $\sum_k a_k e_k$.

Producten tegenover doorsnede: $I_1\cdots I_n \subseteq \bigcap
I_k$ geldt altijd. Omgekeerd volstaat het per inductie het geval $n
= 2$ te behandelen (men gaat na dat $I_1$ en $I_2\cdots I_n$
comaximaal zijn: vermenigvuldigen van $1 = u_l + v_l$ over $l \geq
2$ geeft $1 \in I_1 + I_2\cdots I_n$). Voor $n = 2$: schrijf $1 = u
+ v$ met $u \in I_1$ en $v \in I_2$; voor $x \in I_1 \cap I_2$ is
dan $x = xu + xv \in I_2I_1 + I_1I_2 = I_1I_2$.
\end{proof}

\begin{example}\label{ex:b3:rings:crtz}
In $\Z$ met $I_k = (m_k)$ en de $m_k$ paarsgewijs onderling
ondeelbaar: $\Z/(m_1\cdots m_n)\Z \cong \prod \Z/m_k\Z$ --- de
\omterm{thm:b3:rings:crt}{Chinese reststelling} uit het volume van bachelorjaar 2. Beperkt tot
de eenheden: $(\Z/mn\Z)^\times \cong (\Z/m\Z)^\times \times
(\Z/n\Z)^\times$ voor $\gcd(m,n)=1$, waaruit de multiplicativiteit
van de functie $\varphi$ van Euler volgt (\cref{exo:b3:rings:8}).
\end{example}

\section{Deelbaarheid: euclidisch, hoofdideaal, factorieel}

\begin{definition}\label{def:b3:rings:divisibility}
Zij $A$ een integriteitsdomein en $a, b \in A$. We zeggen dat $a$
het element $b$ \emph{deelt} ($a \mid b$) als $b \in (a) = aA$. Twee
elementen $a, b$ heten \emph{geassocieerd}\index{geassocieerde
elementen} als $a = ub$ met $u \in A^\times$ (gelijkwaardig: $(a) =
(b)$). Een $p \neq 0$ die geen eenheid is, heet:
\begin{itemize}
  \item \emph{irreducibel}\index{irreducibel element} als $p = ab$
        afdwingt dat $a \in A^\times$ of $b \in A^\times$;
  \item \emph{priem}\index{priemelement} als $p \mid ab$ afdwingt
        dat $p \mid a$ of $p \mid b$ (dat wil zeggen: het \omterm{def:b3:rings:ideal}{ideaal}
        $(p)$ is priem).
\end{itemize}
\end{definition}

\begin{proposition}\label{prop:b3:rings:primeirred}
In elk domein geldt: priem $\Rightarrow$ \omterm{def:b3:rings:divisibility}{irreducibel}. Het
omgekeerde is in het algemeen onjuist: in $\Z[\iu\sqrt 5] = \{a +
\iu b\sqrt5 : a, b \in \Z\}$ is $2$ \omterm{def:b3:rings:divisibility}{irreducibel} maar niet priem.
\end{proposition}

\begin{proof}
Zij $p$ priem en $p = ab$. Dan is $p \mid ab$, zeg $p \mid a$: dus
$a = pc$, waaruit $p = pcb$, en wegschrappen van $p$ (het is een
domein!) geeft $cb = 1$: $b \in A^\times$.

In $\Z[\iu\sqrt5]$ gebruiken we de norm $N(x + \iu y\sqrt 5) = x^2
+ 5y^2$, die multiplicatief is (het is $\abs z^2$). Is $2 = ab$ met
$a, b$ geen eenheden, dan is $4 = N(a)N(b)$ met $N(a), N(b) \neq 1$
(de elementen van norm $1$ zijn $\pm1$, de eenheden), dus $N(a) =
2$: onmogelijk, want $x^2 + 5y^2 = 2$ heeft geen gehele
oplossingen. Dus is $2$ \omterm{def:b3:rings:divisibility}{irreducibel}. Maar $2 \mid 6 = (1 +
\iu\sqrt5)(1 - \iu\sqrt5)$ terwijl $2$ geen van beide factoren
deelt ($\frac12 \pm \frac{\iu\sqrt5}2 \notin \Z[\iu\sqrt5]$): niet
priem.
\end{proof}

\begin{definition}\label{def:b3:rings:pidufd}
Een integriteitsdomein $A$ heet:
\begin{itemize}
  \item \emph{euclidisch}\index{euclidisch domein} als er een
        afbeelding $\nu \colon A \setminus \{0\} \to \N$ bestaat
        (een \emph{euclidische functie}) zodanig dat er voor alle
        $a, b$ met $b \ne 0$ een $q$ en een $r$ zijn met $a = bq +
        r$ en ($r = 0$ of $\nu(r) < \nu(b)$);
  \item \emph{hoofdideaaldomein} (een
        \emph{PID}\index{hoofdideaaldomein}) als elk \omterm{def:b3:rings:ideal}{ideaal} van de
        vorm $(a)$ is;
  \item \emph{factorieel} (een \emph{UFD}\index{uniek
        factorisatiedomein}) als elk element $\neq 0$ dat geen
        eenheid is, een product van irreducibele elementen is, op
        de volgorde en op associatie na uniek.
\end{itemize}
\end{definition}

\begin{theorem}\label{thm:b3:rings:euclideanpid}
\omterm{def:b3:rings:pidufd}{Euclidisch} $\Rightarrow$ \omterm{def:b3:rings:pidufd}{hoofdideaaldomein}.
\end{theorem}

\begin{proof}
Zij $I \neq (0)$ een \omterm{def:b3:rings:ideal}{ideaal} en $b \in I \setminus\{0\}$ met
$\nu(b)$ minimaal. Deel voor $a \in I$: $a = bq + r$; dan is $r = a
- bq \in I$, en $\nu(r) < \nu(b)$ zou de minimaliteit
tegenspreken, dus $r = 0$ en $a \in (b)$: $I = (b)$.
\end{proof}

\begin{example}\label{ex:b3:rings:euclidean}
$\Z$ (met $\nu = \abs\cdot$) en $K[X]$ (met $\nu = \deg$) zijn
\omterm{def:b3:rings:pidufd}{euclidisch} --- het volume van bachelorjaar 2 bewees beide delingen.
Ook $\Z[\iu]$ is het, met $\nu = N$ de kwadratische norm
(\cref{exo:b3:rings:4}); de meetkunde achter het bewijs staat in de
figuur hieronder. Een \omterm{def:b3:rings:pidufd}{hoofdideaaldomein} dat niet \omterm{def:b3:rings:pidufd}{euclidisch} is
bestaat, maar is lastig te certificeren (het standaardvoorbeeld is
$\Z\bigl[\frac{1+\iu\sqrt{19}}2\bigr]$); een \omterm{def:b3:rings:pidufd}{UFD} dat geen
\omterm{def:b3:rings:pidufd}{hoofdideaaldomein} is, ligt voor het grijpen: $K[X, Y]$
(\cref{exo:b3:rings:6}) of $\Z[X]$.
\end{example}

\begin{omfigure}
\begin{tikzpicture}[scale=1.05]
  \clip (-2.7,-1.7) rectangle (2.7,2.2);
  \foreach \x in {-2,...,2}
    \foreach \y in {-1,...,2}
      \fill[black!55] (\x,\y) circle (1.4pt);
  \draw[omDef, thick] (0.6,0.4) circle (1);
  \fill[omThm] (0.6,0.4) circle (1.8pt)
    node[below right, font=\small, omThm] {$a/b$};
  \fill[omDef] (1,0) circle (2pt)
    node[below, font=\small, omDef] {$q$};
  \draw[omDef, ->, shorten >=2.5pt, shorten <=2.5pt]
    (0.6,0.4) -- (1,0);
\end{tikzpicture}

{\small Deling bij de gehele getallen van Gauss: het exacte
quotiënt $a/b \in \C$ ligt op afstand $\leq \frac{\sqrt2}{2} < 1$
van een roosterpunt $q \in \Z[\iu]$; dan heeft $r = a - bq$ de norm
$N(r) = N(b)\,\abs{a/b - q}^2 < N(b)$. Eén euclidische deling, en
dus een hele rekenkunde.}
\end{omfigure}

\begin{lemma}[Stijgende ketens van hoofdidealen]
\label{lem:b3:rings:acc}
In een \omterm{def:b3:rings:pidufd}{hoofdideaaldomein} is elke stijgende rij \omterm{def:b3:rings:ideal}{idealen} $I_1
\subseteq I_2 \subseteq \cdots$ vanaf zeker moment constant.
\end{lemma}

\begin{proof}
$I = \bigcup_n I_n$ is een \omterm{def:b3:rings:ideal}{ideaal} (de vereniging is stijgend), dus
$I = (a)$; het element $a$ ligt in een zekere $I_N$, en dan is $I =
(a) \subseteq I_N \subseteq I_n \subseteq I$ voor $n \geq N$.
\end{proof}

\begin{lemma}[Bézout; lemma van Euclides]\label{lem:b3:rings:bezout}
Zij $A$ een \omterm{def:b3:rings:pidufd}{hoofdideaaldomein} en $a, b \in A$. Dan is $(a) + (b) =
(d)$ voor zekere $d$, een \emph{grootste gemene
deler}\index{grootste gemene deler}: $d \mid a$, $d \mid b$, en
elke gemeenschappelijke deler van $a$ en $b$ deelt $d$; bovendien
is $d = au + bv$ voor zekere $u, v$ (Bézout). Bijgevolg is elk
\omterm{def:b3:rings:divisibility}{irreducibel element} van een \omterm{def:b3:rings:pidufd}{hoofdideaaldomein} priem.
\end{lemma}

\begin{proof}
$(a) + (b)$ is een \omterm{def:b3:rings:ideal}{ideaal}, dus van de vorm $(d)$; uit $a, b \in
(d)$ volgt $d \mid a$ en $d \mid b$; en $d = au + bv \in (a) +
(b)$. Een gemeenschappelijke deler $c$ van $a$ en $b$ deelt $au +
bv = d$.

Euclides: zij $p$ \omterm{def:b3:rings:divisibility}{irreducibel} met $p \mid ab$ en $p \nmid a$. Een
\omterm{lem:b3:rings:bezout}{grootste gemene deler} $d$ van $p$ en $a$ deelt $p$, dus is $d$ een
eenheid of \omterm{def:b3:rings:divisibility}{geassocieerd} met $p$ (irreducibiliteit); het tweede is
uitgesloten door $p \nmid a$. Dus $1 = pu + av$, waaruit $b = pub +
abv$, en $p$ deelt beide termen: $p \mid b$.
\end{proof}

\begin{theorem}\label{thm:b3:rings:pidufd}
\omterm{def:b3:rings:pidufd}{Hoofdideaaldomein} $\Rightarrow$ factorieel.
\end{theorem}

\begin{proof}
\emph{Bestaan.} Stel dat een element $a \neq 0$ dat geen eenheid
is, geen ontbinding in irreducibele factoren heeft. Dan is $a$ niet
\omterm{def:b3:rings:divisibility}{irreducibel}: $a = a_1b_1$ met beide factoren geen eenheid; minstens
één ervan, zeg $a_1$, heeft opnieuw geen ontbinding (een product
van twee ontbindbare elementen is ontbindbaar). Herhaald toepassen
levert $a = a_0, a_1, a_2, \dots$, elk een echte deler van de
vorige en zonder ontbinding, dus $(a_0) \subsetneq (a_1) \subsetneq
(a_2) \subsetneq \cdots$ --- de inclusies zijn strikt, want uit
$a_n = a_{n+1}c$ met $c$ geen eenheid zou $(a_n) = (a_{n+1})$
afdwingen dat $c \in A^\times$ (wegschrappen in een domein). Dat
spreekt \cref{lem:b3:rings:acc} tegen.

\emph{Eenduidigheid.} Zij $p_1 \cdots p_r = q_1 \cdots q_s$ met
alle factoren \omterm{def:b3:rings:divisibility}{irreducibel} en $r \leq s$; we gebruiken inductie naar
$r$. Het \omterm{def:b3:rings:divisibility}{priemelement} $p_1$ (\cref{lem:b3:rings:bezout}) deelt het
rechterlid, dus een zekere $q_j$; hernummer zo dat $j = 1$. Omdat
$q_1$ \omterm{def:b3:rings:divisibility}{irreducibel} is en $p_1$ geen eenheid, is $q_1 = u p_1$ met $u
\in A^\times$: $p_1$ en $q_1$ zijn \omterm{def:b3:rings:divisibility}{geassocieerd}. Schrap $p_1$ weg:
$p_2 \cdots p_r = (u q_2) q_3\cdots q_s$, en besluit per inductie
($r = 1$ dwingt $s = 1$ af: een eenheid maal irreducibele elementen
kan geen $1$ zijn).
\end{proof}

\begin{remark}\label{rem:b3:rings:ufdgcd}
In een \omterm{def:b3:rings:pidufd}{UFD} bestaan \omterm{lem:b3:rings:bezout}{grootste gemene delers} (neem de minimale
exponenten in de ontbindingen) en geldt het lemma van Euclides ---
\omterm{def:b3:rings:divisibility}{irreducibel} $=$ priem (\cref{exo:b3:rings:2}) --- maar Bézout kan
falen: in $\Z[X]$ is $\gcd(2, X) = 1$ en toch $1 \neq 2U + XV$
(evalueer in $X = 0$: $1 = 2U(0)$, onmogelijk). Bézout-identiteiten
zijn het exclusieve voorrecht van \omterm{def:b3:rings:pidufd}{hoofdideaaldomeinen}.
\end{remark}

\begin{example}[Een ring zonder eenduidige ontbinding]
\label{ex:b3:rings:nonufd}
Geen van de implicaties \omterm{def:b3:rings:pidufd}{euclidisch} $\Rightarrow$
\omterm{def:b3:rings:pidufd}{hoofdideaaldomein} $\Rightarrow$ \omterm{def:b3:rings:pidufd}{UFD} is een gelijkwaardigheid, en
het falen van de laatste is het waard om één keer helemaal uit
te schrijven. In
\[
A = \Z[\iu\sqrt5] = \{a + \iu b\sqrt5 : a, b \in \Z\},
\qquad N(a + \iu b\sqrt5) = a^2 + 5b^2,
\]
is de norm multiplicatief en geldt $N(z) = 1$ precies wanneer
$z \in A^\times = \{\pm1\}$. Beschouw
\[
6 = 2 \cdot 3 = (1 + \iu\sqrt5)(1 - \iu\sqrt5).
\]
Alle vier de factoren zijn \omterm{def:b3:rings:divisibility}{irreducibel}: hun normen zijn $4, 9,
6, 6$, en een echte ontbinding $z = z_1z_2$ zou $N(z_1) \in
\{2, 3\}$ afdwingen --- maar $a^2 + 5b^2$ is nooit $2$ of $3$
($b = 0$ laat de niet-kwadraten $2, 3$ over; $\abs b \geq 1$
geeft $\geq 5$). En toch is $2$ met geen van beide $1 \pm
\iu\sqrt5$ \omterm{def:b3:rings:divisibility}{geassocieerd} (normen $4 \neq 6$): twee werkelijk
verschillende ontbindingen van $6$ in irreducibele factoren.
Gelijkwaardig: hier is \omterm{def:b3:rings:divisibility}{irreducibel} $\neq$ priem, want $2$ deelt
het product $(1 + \iu\sqrt5)(1 - \iu\sqrt5) = 6$ maar geen van
beide factoren (opnieuw de normen). De ideaaltheoretische
reparatie van dit falen --- \emph{\omterm{def:b3:rings:ideal}{idealen}} ontbinden in plaats
van elementen --- is de geboorte van de algebraïsche
getaltheorie; op ons niveau ijkt het voorbeeld hoe bijzonder de
euclidische ringen $\Z$, $K[X]$ en $\Z[\iu]$ van dit hoofdstuk
werkelijk zijn.
\end{example}

\begin{method}\label{met:b3:rings:quotient}
Om een \omterm{def:b3:rings:ideal}{quotiëntring} $A/I$ te herkennen: zoek een surjectief
morfisme $f \colon A \to B$ met kern $I$ en roep
\cref{thm:b3:rings:firstiso} in; is $A = C[X]$ een veeltermring,
dan is $f$ meestal een evaluatie. Zo is $\Z[X]/(X^2+1) \cong
\Z[\iu]$ (evalueer in $\iu$), $K[X,Y]/(Y - X^2) \cong K[X]$
(evalueer $Y$ in $X^2$) en $\R[X]/(X^2+1) \cong \C$. Om aan te
tonen dat $I$ priem of \omterm{def:b3:rings:primemaximal}{maximaal} is, toon je aan dat het quotiënt
een domein respectievelijk een lichaam is
(\cref{prop:b3:rings:primemaximal}).
\end{method}

\section{Veeltermen over een UFD: Gauss en Eisenstein}

In deze paragraaf is $A$ overal een \omterm{def:b3:rings:pidufd}{UFD} met breukenlichaam $K$
(geconstrueerd als het lichaam van formele quotiënten $a/b$ met $b
\neq 0$, net zoals $\Q$ uit $\Z$; het volume van bachelorjaar 2
voerde die constructie uit voor $\Q$, en ze gaat letterlijk over).
Ons doel: factorialiteit gaat over van $A$ naar $A[X]$, en
irreducibiliteit over $A$ is in wezen irreducibiliteit over het
grotere lichaam $K$.

\begin{definition}\label{def:b3:rings:content}
De \emph{inhoud}\index{inhoud van een veelterm} $c(P)$ van een
$P \in A[X]$ met $P \neq 0$ is een \omterm{lem:b3:rings:bezout}{grootste gemene deler} van haar
coëfficiënten (op een eenheid na bepaald); $P$ heet
\emph{primitief} als $c(P) \in A^\times$. Elke $P \in A[X]$
schrijft zich als $P = c(P)\,P_1$ met $P_1$ primitief, en elke $P
\in K[X]\setminus\{0\}$ als $P = \lambda P_1$ met $\lambda \in
K^\times$ en $P_1 \in A[X]$ primitief (werk de noemers weg en haal
dan de inhoud buiten).
\end{definition}

\begin{lemma}[Gauss]\label{lem:b3:rings:gauss}
Het product van twee primitieve veeltermen van $A[X]$ is
primitief; bijgevolg is $c(PQ) = c(P)c(Q)$ op eenheden
na.\index{lemma van Gauss}
\end{lemma}

\begin{proof}
Zij $P, Q$ primitief en stel dat een irreducibele (dus, in een \omterm{def:b3:rings:pidufd}{UFD},
priem) $p$ alle coëfficiënten van $PQ$ deelt. Reduceer modulo $p$:
in $(A/(p))[X]$ is $\bar P \bar Q = 0$. Maar $A/(p)$ is een domein
($(p)$ is priem), dus is $(A/(p))[X]$ een domein (de
kopcoëfficiënten vermenigvuldigen), zodat $\bar P = 0$ of $\bar Q =
0$: $p$ deelt alle coëfficiënten van $P$ of alle van $Q$, in strijd
met de primitiviteit. Voor het gevolg: schrijf $P = c(P)P_1$ en $Q
= c(Q)Q_1$; dan is $PQ = c(P)c(Q) P_1Q_1$ met $P_1Q_1$ primitief.
\end{proof}

\begin{theorem}\label{thm:b3:rings:gaussufd}
Zij $A$ een \omterm{def:b3:rings:pidufd}{UFD} met breukenlichaam $K$.
\begin{enumerate}
  \item Een primitieve $P \in A[X]$ van graad $\geq 1$ is
        \omterm{def:b3:rings:divisibility}{irreducibel} in $A[X]$ dan en slechts dan als ze
        \omterm{def:b3:rings:divisibility}{irreducibel} is in $K[X]$.
  \item $A[X]$ is een \omterm{def:b3:rings:pidufd}{UFD}; de irreducibele elementen ervan zijn de
        irreducibele elementen van $A$ en de primitieve veeltermen
        die over $K$ \omterm{def:b3:rings:divisibility}{irreducibel} zijn. In het bijzonder zijn
        $\Z[X]$ en, per inductie, $K[X_1, \dots, X_n]$ en $\Z[X_1,
        \dots, X_n]$ \omterm{def:b3:rings:pidufd}{UFD}'s.
\end{enumerate}
\end{theorem}

\begin{proof}
(1) ($\Leftarrow$) Is $P = QR$ in $A[X]$ met $Q, R$ geen eenheden,
dan is geen van beide factoren constant (een constante factor van
een primitieve veelterm is een eenheid), dus is de ontbinding ook
in $K[X]$ echt. ($\Rightarrow$) Stel $P = QR$ met $Q, R \in K[X]$
van graad $\geq 1$. Schrijf $Q = \lambda Q_1$ en $R = \mu R_1$ met
$Q_1, R_1 \in A[X]$ primitief: dan is $P = \lambda\mu\, Q_1R_1$, en
$Q_1R_1$ is primitief volgens Gauss. Nemen we de \omterm{def:b3:rings:content}{inhouden}, dan is
$\lambda\mu \in A^\times$ (beide leden hebben \omterm{def:b3:rings:content}{inhoud} een eenheid;
formeel: $\lambda\mu = c(P) \in A^\times$ op eenheden na, en in het
bijzonder $\lambda \mu \in A$): $P = (\lambda\mu Q_1) R_1$ is dus
een echte ontbinding in $A[X]$.

(2) Bestaan: zij $P \neq 0$ geen eenheid; ontbind $P = c(P)P_1$,
ontbind $c(P)$ in irreducibele elementen van $A$, en ontbind $P_1$
in het \omterm{def:b3:rings:pidufd}{UFD} $K[X]$ als $\prod Q_i$ met $Q_i \in K[X]$ \omterm{def:b3:rings:divisibility}{irreducibel};
schrijven we $Q_i = \lambda_i R_i$ met $R_i \in A[X]$ primitief (en
dus \omterm{def:b3:rings:divisibility}{irreducibel} over $K$, en dus in $A[X]$ volgens (1)), dan is het
product $\prod \lambda_i$ zoals eerder een eenheid van $A$ en is
$P_1 = u\prod R_i$. Eenduidigheid: vergelijk het constante deel en
het veeltermdeel van een ontbinding; de constanten
vermenigvuldigen tot $c(P)$ (Gauss), uniek wegens de
factorialiteit van $A$; de veeltermdelen geven twee ontbindingen
in $K[X]$ van dezelfde veelterm, die dus op constanten van
$K^\times$ na overeenkomen (factorialiteit van $K[X]$,
\cref{thm:b3:rings:pidufd}), en primitieve veeltermen die in
$K[X]$ \omterm{def:b3:rings:divisibility}{geassocieerd} zijn, zijn dat ook in $A[X]$: is $R = \lambda
R'$ met $R, R'$ primitief en $\lambda \in K^\times$, dan dwingt het
nemen van \omterm{def:b3:rings:content}{inhouden} $\lambda \in A^\times$ af.
\end{proof}

\begin{theorem}[Irreducibiliteitscriteria]\label{thm:b3:rings:criteria}
Zij $A$ een \omterm{def:b3:rings:pidufd}{UFD}, $K$ haar breukenlichaam, en $P = a_nX^n + \dots +
a_0 \in A[X]$ primitief van graad $n \geq 1$.
\begin{enumerate}
  \item (\emph{Reductie}) Is $p \in A$ priem met $p \nmid a_n$, en
        is de reductie $\bar P$ \omterm{def:b3:rings:divisibility}{irreducibel} in $(A/(p))[X]$, dan is
        $P$ \omterm{def:b3:rings:divisibility}{irreducibel} in $K[X]$ (en dus in $A[X]$).
  \item (\emph{Eisenstein}\index{criterium van Eisenstein}) Is er
        een priemgetal $p$ met $p \nmid a_n$, $p \mid a_i$ voor $0
        \leq i < n$ en $p^2 \nmid a_0$, dan is $P$ \omterm{def:b3:rings:divisibility}{irreducibel} in
        $K[X]$ (en dus in $A[X]$).
\end{enumerate}
\end{theorem}

\begin{proof}
Volgens \cref{thm:b3:rings:gaussufd}(1) levert een echte
ontbinding over $K$ een $P = QR$ met $Q, R \in A[X]$ en $\deg Q,
\deg R \geq 1$ (constanten zijn uitgesloten: die zouden eenheden
zijn of de primitiviteit bederven).

(1) Reduceer modulo $p$: $\bar P = \bar Q\bar R$ in $(A/(p))[X]$.
Omdat $p \nmid a_n$ en de graad onder reductie alleen kan dalen, is
$\deg \bar Q = \deg Q \geq 1$ en $\deg\bar R = \deg R \geq 1$ (hun
kopcoëfficiënten vermenigvuldigen tot $\bar a_n \neq 0$, dus daalt
geen van beide): $\bar P$ ontbindt echt --- tegenspraak.

(2) Reduceer modulo $p$: $\bar Q \bar R = \bar P = \bar a_n X^n$
(alle lagere coëfficiënten sterven weg). In het domein
$(A/(p))[X]$ zijn de ontbindingen van $cX^n$ ($c \ne 0$) alleen
die in constanten en zuivere machten $c'X^k$: immers, als $\bar
Q\bar R = \bar a_nX^n$ en $\bar Q$ een coëfficiënt $\neq 0$ had in
een graad $< \deg\bar Q$, neem dan de laagste niet-nulle termen:
$\operatorname{val}(\bar Q\bar R) = \operatorname{val}\bar Q +
\operatorname{val}\bar R$ (domein), en dat moet gelijk zijn aan $n
= \deg\bar Q + \deg\bar R$, wat $\operatorname{val} = \deg$
afdwingt voor beide: het zijn allebei monomen. Zoals hierboven
dalen de graden niet, dus zijn de constante termen $Q(0)$ en
$R(0)$ van $Q$ en $R$ door $p$ deelbaar --- allebei, omdat beide
reducties monomen van graad $\geq 1$ zijn. Dan is $p^2 \mid
Q(0)R(0) = a_0$: tegenspraak.
\end{proof}

\begin{example}\label{ex:b3:rings:eisenstein}
$X^n - p$ is \omterm{def:b3:rings:divisibility}{irreducibel} over $\Q$ voor elk priemgetal $p$ en elke
$n \geq 1$ (\omterm{thm:b3:rings:criteria}{Eisenstein} bij $p$): er zijn dus over $\Q$
irreducibele veeltermen van elke graad --- in scherp contrast met
$\C$ (graad $1$, d'Alembert--Gauss, bewezen in
\cref{ch:b3:holomorphic}) en $\R$ (graden $1$ en $2$). De truc van
het \emph{verschuiven} vergroot het bereik van \omterm{thm:b3:rings:criteria}{Eisenstein}: de
$p$-de \emph{cyclotomische veelterm}\index{cyclotomische veelterm}
$\Phi_p = X^{p-1} + \dots + X + 1 = \frac{X^p - 1}{X - 1}$ voldoet
aan
\[
\Phi_p(X + 1) = \frac{(X+1)^p - 1}{X}
= X^{p-1} + \binom{p}{1}X^{p-2} + \dots + \binom{p}{p-1},
\]
en dat is \omterm{thm:b3:rings:criteria}{Eisenstein} bij $p$ ($p \mid \binom pk$ voor $0 < k < p$,
en $\binom{p}{p-1} = p \not\equiv 0 \bmod p^2$): dus is
$\Phi_p(X+1)$, en daarmee $\Phi_p$, \omterm{def:b3:rings:divisibility}{irreducibel} over $\Q$. Dit is
het algebraïsche hart van het verhaal over de regelmatige
$17$-hoek in \cref{ch:b3:galois}.
\end{example}

\begin{method}\label{met:b3:rings:irreducibility}
Om te bewijzen dat $P \in \Z[X]$ \omterm{def:b3:rings:divisibility}{irreducibel} is over $\Q$: (i) maak
$P$ primitief; (ii) probeer \omterm{thm:b3:rings:criteria}{Eisenstein}, op $P(X)$ en op de
verschuivingen $P(X \pm 1)$; (iii) probeer reductie modulo kleine
priemgetallen die de kopcoëfficiënt niet delen ---
irreducibiliteit modulo \emph{één} $p$ volstaat, en over $\mathbb
F_p$ is irreducibiliteit een eindige controle (geen wortels sluit
factoren van graad $1$ uit; test daarna de eindig veel factoren
van elke graad $\leq \deg P/2$); (iv) helpt niets, dan
onbepaalde coëfficiënten. Pas op: reducibiliteit modulo elk
priemgetal impliceert \emph{niet} reducibiliteit over $\Q$
(\cref{exo:b3:rings:11}).
\end{method}

\section{Noetherse ringen}

\begin{definition}\label{def:b3:rings:noetherian}
Een ring $A$ heet \emph{noethers}\index{noetherse ring} als elk
\omterm{def:b3:rings:ideal}{ideaal} van $A$ eindig voortgebracht is.
\end{definition}

\begin{proposition}\label{prop:b3:rings:noethacc}
$A$ is \omterm{def:b3:rings:noetherian}{noethers} dan en slechts dan als elke stijgende rij \omterm{def:b3:rings:ideal}{idealen}
vanaf zeker moment constant is (de \emph{stijgende
ketenvoorwaarde}), dan en slechts dan als elke niet-lege familie
\omterm{def:b3:rings:ideal}{idealen} een \omterm{def:b3:rings:primemaximal}{maximaal} element heeft (voor de inclusie).
\end{proposition}

\begin{proof}
\emph{(Eindig voortgebracht $\Rightarrow$ ketenvoorwaarde)}: voor
een keten $I_1 \subseteq I_2 \subseteq \cdots$ is de vereniging $I$
een \omterm{def:b3:rings:ideal}{ideaal}, voortgebracht door $x_1, \dots, x_r$; alle $x_i$ liggen
in een zekere $I_N$, dus $I = I_N = I_n$ voor $n \geq N$.
\emph{(Ketenvoorwaarde $\Rightarrow$ maximale elementen)}: had een
niet-lege familie $\mathcal F$ geen \omterm{def:b3:rings:primemaximal}{maximaal} element, kies dan $I_1
\in \mathcal F$ en inductief $I_{n+1} \supsetneq I_n$ in $\mathcal
F$ (mogelijk omdat $I_n$ niet \omterm{def:b3:rings:primemaximal}{maximaal} is): een oneindige strikt
stijgende keten. (Hier gebruiken we het axioma van afhankelijke
keuzen, een zwakke vorm van het keuzeaxioma waar we niet over
tobben.) \emph{(Maximale elementen $\Rightarrow$ eindig
voortgebracht)}: zij $I$ een \omterm{def:b3:rings:ideal}{ideaal}; de familie van eindig
voortgebrachte \omterm{def:b3:rings:ideal}{idealen} binnen $I$ is niet leeg ($(0)$), en een
\omterm{def:b3:rings:primemaximal}{maximaal} element $J = (x_1, \dots, x_r)$ moet gelijk zijn aan $I$:
anders levert het toevoegen van een $x \in I \setminus J$ aan de
voortbrengers een strikt groter lid van de familie.
\end{proof}

\begin{theorem}[Basisstelling van Hilbert]\label{thm:b3:rings:hilbert}
Is $A$ \omterm{def:b3:rings:noetherian}{noethers}, dan ook $A[X]$. Bijgevolg zijn ook $A[X_1, \dots,
X_n]$ en al hun quotiënten \omterm{def:b3:rings:noetherian}{noethers}.\index{basisstelling van
Hilbert}
\end{theorem}

\begin{proof}
Zij $I$ een \omterm{def:b3:rings:ideal}{ideaal} van $A[X]$ en stel dat $I$ niet eindig
voortgebracht is. Bouw een rij: neem $f_1 \in I \setminus \{0\}$
van minimale graad, en inductief $f_{k+1} \in I \setminus (f_1,
\dots, f_k)$ van minimale graad (die verzameling is per
veronderstelling niet leeg). De graden $d_k = \deg f_k$ zijn
niet-dalend: immers $f_{k+1} \notin (f_1,\dots,f_k) \supseteq (f_1,
\dots, f_{k-1})$, dus deed $f_{k+1}$ in stap $k$ al mee en verloor
of speelde gelijk: $d_{k+1} \geq d_k$. Zij $a_k \in A$ de
kopcoëfficiënt van $f_k$. De keten \omterm{def:b3:rings:ideal}{idealen} $(a_1) \subseteq (a_1,
a_2) \subseteq \cdots$ stabiliseert: voor zekere $n$ is $a_{n+1}
\in (a_1, \dots, a_n)$, zeg $a_{n+1} = \sum_{k\leq n} u_k a_k$.
Beschouw
\[
g = f_{n+1} - \sum_{k=1}^{n} u_k X^{\,d_{n+1} - d_k} f_k .
\]
Dan is $g \in I \setminus (f_1, \dots, f_n)$ (de som ligt in het
\omterm{def:b3:rings:ideal}{ideaal}, $f_{n+1}$ niet), en toch valt de coëfficiënt van graad
$d_{n+1}$ weg: $\deg g < d_{n+1}$, in strijd met de minimaliteit
van $d_{n+1} = \deg f_{n+1}$.

Herhaald toepassen geeft dat $A[X_1, \dots, X_n] = (A[X_1, \dots,
X_{n-1}])[X_n]$ \omterm{def:b3:rings:noetherian}{noethers} is; en een quotiënt $A/I$ is \omterm{def:b3:rings:noetherian}{noethers}
omdat zijn \omterm{def:b3:rings:ideal}{idealen} $J/I$ zich optillen tot \omterm{def:b3:rings:ideal}{idealen} van $A$
(correspondentie), waar eindig veel voortbrengers op voortbrengers
projecteren.
\end{proof}

\begin{remark}\label{rem:b3:rings:noethmeaning}
Noetheriaanheid is het eindigheidsaxioma van de algebraïsche
meetkunde: elk stelsel veeltermvergelijkingen in $n$
veranderlijken, hoe oneindig ook, is gelijkwaardig met eindig veel
ervan --- de oplossingsverzameling wordt door eindig veel
veeltermen uitgesneden. \omterm{def:b3:rings:pidufd}{Hoofdideaaldomeinen} zijn \omterm{def:b3:rings:noetherian}{noethers}
(triviaal); $\Z[X_1, X_2, \dots]$ in oneindig veel veranderlijken
is dat niet ($(X_1) \subsetneq (X_1, X_2) \subsetneq \cdots$). Ook
in de analyse komen niet-noetherse ringen vanzelf voor: de continue
functies op $\intcc01$ vormen er een (\cref{exo:b3:rings:10}).
\end{remark}

\section{Oefeningen}

\begin{exercise}[$\star$]\label{exo:b3:rings:1}
Herken de quotiënten: (a) $\Z[X]/(X^2 + 1) \cong \Z[\iu]$;
(b) $\R[X]/(X^2+1) \cong \C$; (c) $\mathbb F_2[X]/(X^2 + X + 1)$ is
een lichaam met $4$ elementen --- schrijf de
vermenigvuldigingstabel op.
\end{exercise}

\begin{exercise}[$\star$]\label{exo:b3:rings:2}
(a) Toon aan dat in een \omterm{def:b3:rings:pidufd}{UFD} elk \omterm{def:b3:rings:divisibility}{irreducibel element} priem is.
(b) Toon aan dat een eindig integriteitsdomein een lichaam is.
(c) Leid af dat in een eindige ring elk \omterm{def:b3:rings:primemaximal}{priemideaal} \omterm{def:b3:rings:primemaximal}{maximaal} is.
\end{exercise}

\begin{exercise}[$\star$]\label{exo:b3:rings:3}
Ga in $\Z[\iu\sqrt5]$ na dat $3$, $1 + \iu\sqrt5$ en $1 -
\iu\sqrt5$ \omterm{def:b3:rings:divisibility}{irreducibel} zijn en dat $9 = 3\cdot 3 = (2 +
\iu\sqrt5)(2 - \iu\sqrt5)$, en besluit opnieuw (na
\cref{prop:b3:rings:primeirred}) dat $\Z[\iu\sqrt5]$ geen \omterm{def:b3:rings:pidufd}{UFD} is.
Waar precies faalt de eenduidigheid?
\end{exercise}

\begin{exercise}[$\star\star$]\label{exo:b3:rings:4}
(a) Toon aan dat $\Z[\iu]$ \omterm{def:b3:rings:pidufd}{euclidisch} is voor de norm $N(x + \iu y)
= x^2 + y^2$: kies bij gegeven $a, b \neq 0$ een $q \in \Z[\iu]$
die het dichtst bij $a/b \in \C$ ligt.
(b) Bepaal $\Z[\iu]^\times$.
(c) Dezelfde vragen voor $\Z[\iu\sqrt2]$ en $N(x + \iu y\sqrt2) =
x^2 + 2y^2$. Waarom faalt hetzelfde argument voor
$\Z[\iu\sqrt5]$?
\end{exercise}

\begin{exercise}[$\star\star$]\label{exo:b3:rings:5}
Zij $A$ een ring. (a) Toon aan dat $1 + x \in A^\times$ zodra $x$
nilpotent is ($x^n = 0$ voor zekere $n$). (b) Toon aan dat
$A[X]^\times = A^\times$ als $A$ een domein is; geef een
tegenvoorbeeld over $\Z/4\Z$. (c) Toon aan dat een domein geen
andere idempotenten ($e^2 = e$) heeft dan $0$ en $1$, en geen
andere nilpotente elementen dan $0$.
\end{exercise}

\begin{exercise}[$\star\star$]\label{exo:b3:rings:6}
In $A = K[X, Y]$: (a) toon aan dat het \omterm{def:b3:rings:ideal}{ideaal} $(X, Y)$ \omterm{def:b3:rings:primemaximal}{maximaal}
maar geen hoofdideaal is --- zodat $K[X,Y]$ een \omterm{def:b3:rings:pidufd}{UFD}
(\cref{thm:b3:rings:gaussufd}) is dat geen \omterm{def:b3:rings:pidufd}{hoofdideaaldomein} is;
(b) herken $K[X, Y]/(Y - X^2)$ en $K[X,Y]/(XY - 1)$ als deelringen
van rationale functies; (c) is $(Y - X^2)$ priem? \omterm{def:b3:rings:primemaximal}{maximaal}?
\end{exercise}

\begin{exercise}[$\star\star$]\label{exo:b3:rings:7}
\omterm{def:b3:rings:divisibility}{Irreducibel} over $\Q$ of niet: $X^5 - 12X^3 + 36X - 12$;\quad
$X^4 + X + 1$ \emph{(reduceer modulo $2$)};\quad $X^4 + 4$;\quad
$\Phi_8 = X^4 + 1$ \emph{(verschuif over $1$)};\quad $X^3 - X - 1$.
\end{exercise}

\begin{exercise}[$\star\star$]\label{exo:b3:rings:8}
(a) Bewijs met \cref{thm:b3:rings:crt} dat de functie van Euler
multiplicatief is op onderling ondeelbare argumenten en dat
$\varphi(p^k) = p^{k-1}(p-1)$; vind zo $\varphi(n) = n\prod_{p \mid
n}(1 - \frac1p)$ terug.
(b) Los op: $x \equiv 2 \pmod 7$, $x \equiv 5 \pmod{11}$, $x \equiv
1 \pmod{13}$, en geef daarbij de idempotenten $e_k$ uit het bewijs
van \cref{thm:b3:rings:crt}.
\end{exercise}

\begin{exercise}[$\star\star\star$]\label{exo:b3:rings:9}
(Het nilradicaal) Zij $\operatorname{Nil}(A)$ de verzameling
nilpotente elementen. (a) Toon aan dat $\operatorname{Nil}(A)$ een
\omterm{def:b3:rings:ideal}{ideaal} is dat in elk \omterm{def:b3:rings:primemaximal}{priemideaal} ligt. (b) Zij omgekeerd $a$ niet
nilpotent; produceer met het lemma van Zorn, toegepast op de
\omterm{def:b3:rings:ideal}{idealen} die $S = \{a^n : n \in \N\}$ mijden, een \omterm{def:b3:rings:primemaximal}{priemideaal} dat
$a$ niet bevat. Besluit:
\[
\operatorname{Nil}(A) = \bigcap_{\mathfrak p \text{ priem}}
\mathfrak p .
\]
\end{exercise}

\begin{exercise}[$\star\star\star$]\label{exo:b3:rings:10}
(a) Zij $A$ \omterm{def:b3:rings:noetherian}{noethers} en $f \colon A \to A$ een surjectief
ringmorfisme. Toon aan dat $f$ injectief is. \emph{(Beschouw $\ker
f \subseteq \ker f^2 \subseteq \cdots$.)}
(b) Toon aan dat de ring $\mathcal C(\intcc01, \R)$ van continue
functies niet \omterm{def:b3:rings:noetherian}{noethers} is. \emph{(Beschouw $I_n = \{f : f = 0
\text{ op } \intcc0{1/n}\}$.)}
(c) Toon aan dat in een \omterm{def:b3:rings:noetherian}{noethers} \emph{domein} elk element $\neq 0$
dat geen eenheid is, een (eindig) product van irreducibele
elementen is --- zodat de niet-factorialiteit van
$\Z[\iu\sqrt 5]$ uitsluitend een falen van de
\emph{eenduidigheid} is.
\end{exercise}

\begin{exercise}[$\star\star\star$]\label{exo:b3:rings:11}
Zij $P = X^4 + 1$.
(a) Toon aan dat $P$ \omterm{def:b3:rings:divisibility}{irreducibel} is over $\Q$
(\cref{exo:b3:rings:7}).
(b) Toon aan dat $P$ modulo \emph{elk} priemgetal $p$ reducibel is:
behandel $p = 2$; toon vervolgens voor oneven $p$ aan dat $8 \mid
p^2 - 1$ en neem voorlopig zonder bewijs aan (het wordt bewezen in
\cref{ch:b3:galois}) dat de multiplicatieve groep van het lichaam
met $p^2$ elementen cyclisch is, om te besluiten dat $P$ modulo $p$
in twee kwadratische factoren uiteenvalt; maak die expliciet
wanneer een van $-1$, $2$, $-2$ een kwadraat modulo $p$ is, en toon
aan dat dat er altijd een is.
\end{exercise}

\begin{exercise}[$\star\star$]\label{exo:b3:rings:12}
(Idempotenten splitsen ringen) Een element $e$ van een
commutatieve ring $A$ heet \emph{idempotent} als $e^2 = e$.
(a) Toon aan dat met $e$ ook $1 - e$ idempotent is, en dat de
afbeelding $x \mapsto (ex, (1-e)x)$ een ringisomorfisme $A
\cong Ae \times A(1-e)$ is, waarbij $Ae$ een ring met eenheid $e$
is.
(b) Bepaal alle idempotenten van een domein, en die van
$\Z/12\Z$; geef het isomorfisme $\Z/12\Z \cong \Z/4\Z \times
\Z/3\Z$ door zijn twee niet-triviale idempotenten te noemen.
(c) Toon aan dat de ontbinding van $\Z/n\Z$ volgens de \omterm{thm:b3:rings:crt}{Chinese
reststelling} (\cref{ex:b3:rings:crtz}) precies overeenkomt met
de idempotenten $e_i \equiv 1 \bmod p_i^{a_i}$ en $e_i \equiv 0$
modulo de andere priemmachten: ringen vallen uiteen langs hun
idempotenten zoals ruimten uiteenvallen langs projecties.
\end{exercise}

\section{Probleem: de tweekwadratenstelling van Fermat}

\begin{problem}[{Weekendopgave --- sommen van twee kwadraten, via
$\Z[\iu]$}]\label{pb:b3:rings:1}
Welke gehele getallen zijn sommen van twee kwadraten? Het antwoord
van Fermat (1640) is een van de juwelen van de rekenkunde; de
gehele getallen van Gauss maken van het bewijs ringtheorie. Overal
is $N(x + \iu y) = x^2 + y^2$ de norm, is $\Z[\iu]$ \omterm{def:b3:rings:pidufd}{euclidisch}
(\cref{exo:b3:rings:4}) en dus een \omterm{def:b3:rings:pidufd}{hoofdideaaldomein} en een \omterm{def:b3:rings:pidufd}{UFD}, en
betekent \emph{Gauss-priemgetal} een \omterm{def:b3:rings:divisibility}{priemelement} (=\ \omterm{def:b3:rings:divisibility}{irreducibel
element}) van $\Z[\iu]$.

\medskip
\noindent\textbf{Deel I --- Normen en Gauss-priemgetallen.}
\begin{enumerate}
  \item Ga na dat $N(zw) = N(z)N(w)$, leid opnieuw af dat
        $\Z[\iu]^\times = \{\pm1, \pm\iu\}$, en bewijs de
        \emph{identiteit van Brahmagupta}: een product van twee
        sommen van twee kwadraten is zelf een som van twee
        kwadraten.
  \item Toon aan dat $z$ een Gauss-priemgetal is zodra $N(z)$ een
        priemgetal is.
  \item Toon aan dat elk Gauss-priemgetal $\pi$ precies één
        priemgetal $p$ deelt \emph{(beschouw $N(\pi) =
        \pi\bar\pi$)}, en dat dan $N(\pi) \in \{p, p^2\}$.
  \item Leid de tweedeling af: voor elk priemgetal $p$ blijft
        ofwel $p$ priem in $\Z[\iu]$ (en bestaat er geen
        Gauss-priemgetal van norm $p$), ofwel is $p = \pi\bar\pi$
        met $\pi$ een Gauss-priemgetal van norm $p$ --- en dan is
        $p = a^2 + b^2$.
\end{enumerate}

\medskip
\noindent\textbf{Deel II --- De stelling van Wilson en $-1$ modulo
$p$.}
\begin{enumerate}[resume]
  \item Bewijs de \emph{stelling van Wilson}: voor $p$ priem is
        $(p-1)! \equiv -1 \pmod p$. \emph{(Paar elke rest met haar
        inverse; welke zijn met zichzelf gepaard?)}
  \item Zij $p$ een oneven priemgetal en $m = \frac{p-1}2$. Toon
        aan dat $(m!)^2 \equiv (-1)^{m+1} \pmod p$ \emph{(vervang
        in $(p-1)!$ elke factor $k > m$ door $-(p - k)$)}.
  \item Besluit: $-1$ is een kwadraat modulo $p$ dan en slechts dan
        als $p = 2$ of $p \equiv 1 \pmod 4$. \emph{(Voor ``slechts
        dan'': is $x^2 \equiv -1$, wat is dan de orde van $x$ in
        $(\Z/p\Z)^\times$, en wat zegt Lagrange?)}
\end{enumerate}

\medskip
\noindent\textbf{Deel III --- De splitsingswet.}
\begin{enumerate}[resume]
  \item Zij $p \equiv 1 \pmod 4$ en $x$ zodanig dat $p \mid x^2 + 1
        = (x + \iu)(x - \iu)$. Toon aan dat $p$ \emph{geen}
        Gauss-priemgetal is, en besluit met Deel I: $p = a^2 +
        b^2$.
  \item Zij $p \equiv 3 \pmod 4$. Toon rechtstreeks aan dat $p$
        geen som van twee kwadraten is \emph{(kwadraten modulo
        $4$)}, en leid af dat $p$ een Gauss-priemgetal blijft.
  \item Handel $p = 2$ af: geef de ontbinding $2 = -\iu(1+\iu)^2$
        en ga na dat $1 + \iu$ een Gauss-priemgetal is. ($2$ is het
        enige \emph{vertakte} priemgetal: het is op een eenheid na
        deelbaar door het kwadraat van een Gauss-priemgetal.)
  \item Stel de classificatie van de Gauss-priemgetallen samen, op
        eenheden na: $1 + \iu$; de gehele getallen $p \equiv 3
        \pmod 4$; en de toegevoegde paren $\pi, \bar\pi$ van norm
        $p \equiv 1 \pmod 4$. Controleer haar op $5 =
        (2+\iu)(2-\iu)$ en op $3$.
\end{enumerate}

\medskip
\noindent\textbf{Deel IV --- De tweekwadratenstelling.}
\begin{enumerate}[resume]
  \item Bewijs de rechtstreekse helft: verschijnt in de ontbinding
        $n = \prod p_i^{\alpha_i}$ elk priemgetal $\equiv 3 \pmod
        4$ met een even exponent, dan is $n$ een som van twee
        kwadraten. \emph{(Brahmagupta samen met de Delen
        II--III.)}
  \item Bewijs de omkering: is $n = a^2 + b^2 = N(a + \iu b)$ en
        deelt $q \equiv 3 \pmod 4$ het getal $n$, toon dan aan dat
        $q$, een Gauss-priemgetal, $a + \iu b$ of $a - \iu b$
        deelt, dat het in feite zowel $a$ als $b$ deelt, en besluit
        met inductie naar $n$ dat de exponent van $q$ in $n$ even
        is.
  \item Formuleer de uiteindelijke stelling. Welke van $2025$,
        $2026$, $2027$ zijn sommen van twee kwadraten? \emph{($2025
        = 81 \cdot 25$; $2026 = 2 \cdot 1013$ met $1013$ priem;
        $2027$ priem.)}
  \item (Slotwoord) Toon aan dat een priemgetal $p \equiv 1 \pmod
        4$ in wezen op één manier een som van twee kwadraten is:
        is $p = a^2 + b^2 = c^2 + d^2$ met positieve gehele
        getallen, dan is $\{a, b\} = \{c, d\}$. \emph{(De
        eenduidigheid van de ontbinding in $\Z[\iu]$.)}
\end{enumerate}

\medskip
\noindent\textbf{Deel V --- Representaties tellen: de formule van
Jacobi en de reeks van Leibniz.} Schrijf $r_2(n) = \#\{(a, b) \in
\Z^2 : a^2 + b^2 = n\}$ (geordende paren, tekens en nullen
meegerekend), en zij $\chi$ het niet-triviale karakter modulo $4$:
$\chi(d) = +1$ als $d \equiv 1$, $-1$ als $d \equiv 3 \pmod4$, en
$0$ als $d$ even is.
\begin{enumerate}[resume]
  \item (Opwarmer, ter contrast) Welke gehele getallen zijn
        \emph{verschillen} van twee kwadraten? Toon aan: $n = a^2 -
        b^2$ met $a, b \in \Z$ dan en slechts dan als $n
        \not\equiv 2 \pmod 4$ --- geen ringtheorie nodig, en geen
        structuur die te vergelijken is met wat volgt.
  \item Toon aan dat $r_2(n)$ het aantal $z \in \Z[\iu]$ met $N(z)
        = n$ is. Schrijf $n = 2^{a}\prod_jp_j^{b_j}
        \prod_kq_k^{c_k}$ met $p_j \equiv 1$ en $q_k \equiv 3
        \pmod4$, en toon met de classificatie van vraag 11 en de
        eenduidige ontbinding aan dat zulke $z$ bestaan dan en
        slechts dan als alle $c_k$ even zijn, en dat dan
        \[
        r_2(n) = 4\prod_j\,(b_j + 1) .
        \]
        \emph{(Tel: $z = u\,(1+\iu)^{a}\prod_j\pi_j^{s_j}
        \bar\pi_j^{\,b_j - s_j}\prod_kq_k^{c_k/2}$ met $u$ een
        eenheid en $0 \leq s_j \leq b_j$; waarom is die lijst
        volledig en zonder herhaling?)}
  \item Toon aan dat $d \mapsto \chi(d)$ volledig multiplicatief
        is, leid af dat $n \mapsto \sum_{d \mid n}\chi(d)$
        multiplicatief is, en bereken die functie op priemmachten:
        ze is $1$ op $2^a$; $b + 1$ op $p^b$ ($p \equiv 1$); en $1$
        of $0$ op $q^c$ ($q \equiv 3$), naargelang $c$ even of
        oneven is.
  \item Besluit de \emph{stelling van Jacobi}:
        \[
        r_2(n) = 4\sum_{d \mid n}\chi(d) =
        4\bigl(d_1(n) - d_3(n)\bigr),
        \]
        waarbij $d_i(n)$ de delers $\equiv i \pmod 4$ telt.
        Controleer dit voor $n = 3, 5, 9, 25$, en geef de $16$
        representaties van $65$.
  \item (De cirkel) Toon aan dat $\sum_{n \leq x}r_2(n)$ het aantal
        roosterpunten van $\Z^2$ in de gesloten schijf met straal
        $\sqrt x$ is, en bewijs
        \[
        \sum_{n\leq x}r_2(n) = \pi x + O(\sqrt x)
        \]
        \emph{(elk roosterpunt bezit een eenheidsvierkant;
        vergelijk oppervlakten, waarbij de fout in een ring van
        breedte $O(1)$ leeft)}.
  \item (Leibniz, rekenkundig gelezen) Combineer de vragen 19--20:
        \[
        \sum_{d \leq x}\chi(d)\Bigl\lfloor\frac
        xd\Bigr\rfloor = \frac{\pi x}4 + O(\sqrt x),
        \]
        en leid daaruit --- door de gehele delen zorgvuldig weg te
        werken --- de reeks van Leibniz af:
        \[
        1 - \frac13 + \frac15 - \frac17 + \dots = \frac\pi4 .
        \]
        De alternerende reeks van de omgekeerden van de oneven
        getallen \emph{is} het gemiddelde overschot van de delers
        $\equiv 1$ op de delers $\equiv 3$: analyse berekend door
        rekenkunde.
  \item (Hoe zeldzaam zijn sommen van twee kwadraten?) Toon aan dat
        geen enkel geheel getal $\equiv 3 \pmod 4$ een som van twee
        kwadraten is (op twee manieren: kwadraten modulo $4$, of
        het pariteitscriterium van vraag 17), zodat minstens een
        kwart van alle gehele getallen buiten de boot valt; en toon
        aan dat het gemiddelde $\frac1x\sum_{n\leq x}r_2(n) \to
        \pi$ uit vraag 20 verenigbaar is met dichtheid $0$ voor de
        representeerbare getallen --- geef gehele getallen met
        abnormaal veel representaties (neem producten van veel
        priemgetallen $\equiv 1 \bmod 4$) om uit te leggen hoe een
        verdwijnend aandeel toch een positief gemiddelde kan
        dragen. (Landau bewees dat de werkelijke dichtheid als
        $1/\sqrt{\log x}$ afneemt; dat gaat ons gereedschap te
        boven, maar het mechanisme is nu zichtbaar.)
\end{enumerate}

\medskip
\noindent\textbf{Deel VI --- Aanvullingen: primitieve
representaties en Pythagoras.}
\begin{enumerate}[resume]
  \item Noem een representatie $n = a^2 + b^2$ \emph{primitief} als
        $\gcd(a, b) = 1$. Toon aan dat $n \geq 1$ een primitieve
        representatie toelaat dan en slechts dan als $4 \nmid n$ en
        geen priemgetal $q \equiv 3 \pmod 4$ het getal $n$ deelt.
        \emph{(Voor de noodzakelijkheid: hergebruik de afdaling van
        vraag 13 en de kwadraten modulo $4$; voor de
        toereikendheid: bouw $z$ uitsluitend uit $1 + \iu$ en de
        $\pi_j$ --- geen toegevoegden --- en leg uit waarom een
        gemeenschappelijke priemfactor van $a$ en $b$ zowel $\pi_j$
        als $\bar\pi_j$, of $(1+\iu)^2$, in $z$ zou afdwingen.)}
  \item (Pythagoreïsche drietallen) Zij $a^2 + b^2 = c^2$ met $a,
        b, c$ positief, $\gcd(a, b) = 1$ en $b$ even. Toon aan dat
        $a + \iu b$ en $a - \iu b$ onderling ondeelbaar zijn in
        $\Z[\iu]$ \emph{(een gemeenschappelijke priemdeler in
        $\Z[\iu]$ zou $2a$ en $2b$ delen, en $c$ is oneven)}, leid
        uit de eenduidige ontbinding af dat $a + \iu b = u(m + \iu
        n)^2$ voor een eenheid $u$, en besluit met de klassieke
        parametrisatie: op verwisseling van $a$ en $b$ na is
        \[
        a = m^2 - n^2, \qquad b = 2mn, \qquad c = m^2 + n^2,
        \]
        met $m > n \geq 1$ onderling ondeelbaar en van
        tegengestelde pariteit. Vind $(3, 4, 5)$ en $(21, 20, 29)$
        terug uit $(m, n) = (2, 1)$ en $(5, 2)$.
  \item (Numerieke controle) Neem $x = 25$. Bereken $r_2(n)$ voor
        $1 \leq n \leq 25$ met de formule van Jacobi, ga na dat de
        waarden $\neq 0$ precies optreden bij $n = 1, 2, 4, 5, 8,
        9, 10, 13, 16, 17, 18, 20, 25$, en dat
        \[
        \sum_{n \leq 25} r_2(n) = 80
        = 4\sum_{d \leq 25}\chi(d)
        \Bigl\lfloor\frac{25}d\Bigr\rfloor .
        \]
        Controleer dat de gesloten schijf met straal $5$ precies
        $81$ roosterpunten bevat, en vergelijk met $\pi x \approx
        78.5$: de fout blijft ruim binnen de $O(\sqrt x)$ van vraag
        20.
\end{enumerate}
\end{problem}
